\documentclass[12pt,letterpaper]{article}
\usepackage[utf8]{inputenc}
\usepackage[spanish]{babel}
\usepackage{times}
\usepackage[margin=2.54cm]{geometry}
\usepackage{setspace}
\usepackage{titlesec}
\usepackage{enumitem}
\usepackage{hyperref}
\usepackage{xurl}
\usepackage{fancyhdr}
\usepackage{microtype}
\sloppy

\setlength{\parindent}{0pt}
\setlength{\parskip}{6pt}
\singlespacing

\pagestyle{fancy}
\fancyhf{}
\fancyfoot[C]{\thepage}
\renewcommand{\headrulewidth}{0pt}

\titleformat{\section}{\bfseries}{}{0em}{}
\titlespacing{\section}{0pt}{6pt}{2pt}

\hypersetup{colorlinks=true, urlcolor=blue, linkcolor=black}

\begin{document}

\begin{center}
{\bfseries TRES PALABRAS, UN SOLO CANDADO}\\[4pt]
\textit{K.V. Ramírez Ambrosio}\\
7690-22-2239 \quad Universidad Mariano Gálvez\\
Seminario de Tecnologías de Información\\
kramireza10@miumg.edu.gt
\end{center}

\vspace{6pt}

\section*{Resumen}
Este artículo tiene como objetivo aclarar la diferencia entre tres términos que en la práctica suelen usarse como si fueran lo mismo: ciberseguridad, seguridad informática y seguridad de la información. La metodología aplicada consistió en la revisión de fuentes institucionales reconocidas, como el Instituto NIST. Como resultado se identificó que la seguridad de la información es el concepto más amplio, porque protege la información sin importar el formato en que exista, ya sea digital, en papel o incluso hablada; la seguridad informática se enfoca específicamente en proteger los sistemas, equipos y redes donde esa información se almacena o se procesa; y la ciberseguridad se concentra en proteger la información digital que viaja o vive dentro de sistemas interconectados a internet, defendiéndola de ataques provenientes del ciberespacio. Se concluye que estos tres conceptos no son sinónimos, sino que se relacionan como círculos que se traslapan entre sí, y que confundirlos puede llevar a que una organización piense que está protegida en todos los frentes cuando en realidad solo cubrió uno de ellos.

\textbf{Palabras clave:} ciberseguridad, seguridad informática, seguridad de la información, protección de datos, amenazas digitales.

\section*{Desarrollo del tema}

Es muy común escuchar a alguien decir ciberseguridad para referirse a cualquier cosa relacionada con proteger datos, sin darse cuenta de que en realidad está mezclando tres conceptos distintos.Para empezar por el concepto más amplio, la seguridad de la información se encarga de proteger la información en cualquier forma en la que esta exista, no solo la digital. Esto incluye documentos impresos, conversaciones habladas, información almacenada en la nube o incluso datos escritos a mano por ejemplo La Magna Carta. Su objetivo principal es mantener tres características de esa información: que sea confidencial (solo la vea quien debe verla), que sea íntegra (que no se altere sin autorización) y que esté disponible cuando se necesite. Como su alcance es tan amplio, la seguridad de la información suele apoyarse en normas específicas, como la ISO 27001, para organizar todas las políticas y controles necesarios (Hillstone, 2023).

Dentro de ese concepto más amplio aparece la seguridad informática, que ya no se preocupa por toda la información en general, sino específicamente por los sistemas donde esa información vive: computadoras, servidores, software y redes. Su función es proteger el hardware y el software de daños, accesos no autorizados o fallas que puedan comprometer la información que esos sistemas almacenan o procesan (TI Rescue, 2024). Aquí es donde entran temas como los antivirus, los respaldos de información, las contraseñas de acceso a los equipos y el mantenimiento de los sistemas.

Finalmente está la ciberseguridad, que suele confundirse con la seguridad informática porque ambas trabajan muy de cerca, pero no son lo mismo. Según una definición reconocida del NIST, la ciberseguridad consiste en prevenir el daño y proteger los sistemas electrónicos, servicios de comunicación y la información que estos contienen, para garantizar su disponibilidad, integridad y confidencialidad frente a amenazas del entorno digital (Computer Security Resource Center, s.f.). La diferencia clave está en el origen de la amenaza: mientras la seguridad informática protege el sistema en general (incluyendo fallas internas, errores humanos o daños físicos), la ciberseguridad se enfoca puntualmente en amenazas que vienen del ciberespacio, como los ataques de hackers, el malware o el phishing (Telefónica, 2024).

Otra diferencia que vale la pena mencionar es el rol de las personas dentro de cada disciplina. La seguridad informática y la seguridad de la información pueden ser responsabilidad casi exclusiva del departamento de tecnología, pero la ciberseguridad necesita la participación de todos los colaboradores de una empresa, porque un solo clic en un correo falso puede abrir la puerta a un ataque, sin importar qué tan buenos sean los sistemas técnicos de defensa (Hillstone, 2023). Por eso las capacitaciones sobre cómo identificar correos sospechosos son parte esencial de cualquier estrategia de ciberseguridad, y no solo un tema técnico de servidores o software. Esto explica por qué un ataque de phishing es un problema de ciberseguridad, pero también es, al mismo tiempo, un problema de seguridad de la información, porque al final lo que está en riesgo es la información de la organización.

\section*{Observaciones y comentarios}

Algo interesante al revisar este tema es que ni siquiera los expertos se ponen totalmente de acuerdo en los límites exactos entre estos tres conceptos. Algunas fuentes consideran que la ciberseguridad es una parte de la seguridad informática, mientras que otras la presentan casi como sinónimo, y otras más la ubican directamente dentro de la seguridad de la información. Esto sugiere que, más allá de memorizar una definición exacta, lo verdaderamente útil es entender la idea general: que hay diferentes niveles de protección según qué se protege (la información en general, los sistemas, o específicamente el entorno digital conectado a internet).

También llama la atención que buena parte de la confusión entre estos términos ocurre porque en el lenguaje cotidiano, sobre todo en redes sociales y noticias, la palabra ciberseguridad se usa como una especie de paraguas para todo, incluso para temas que en realidad pertenecen a la seguridad informática tradicional, como perder información por una falla de hardware que no tiene nada que ver con un ataque externo.

Como limitación de este trabajo, no se profundizó en marcos normativos específicos como la ISO 27001 o el marco de ciberseguridad del NIST, ya que el objetivo era diferenciar los conceptos y no explicar en detalle cada estándar. Como línea de investigación futura, sería interesante analizar cómo estas tres disciplinas se distribuyen dentro de un organigrama real de una empresa, y quién es responsable de cada una en la práctica.

\section*{Conclusiones}

\begin{enumerate}[label=\arabic*., leftmargin=*]
\item La seguridad de la información es el concepto más amplio de los tres, porque protege la información en cualquier formato en el que exista, sea digital, físico o verbal.
\item La seguridad informática se enfoca específicamente en proteger los sistemas, equipos y redes donde se almacena o procesa la información, sin importar si están conectados a internet o no.
\item La ciberseguridad se concentra en proteger la información digital y los sistemas interconectados frente a amenazas que provienen específicamente del ciberespacio, como ataques de hackers o malware.
\item Los tres conceptos no son sinónimos, sino que se relacionan entre sí como círculos que se traslapan, y confundirlos puede hacer que una organización deje huecos de protección sin cubrir.
\item La ciberseguridad, a diferencia de las otras dos, requiere la participación activa de todos los colaboradores de una organización, porque muchas amenazas digitales dependen del comportamiento humano y no solo de la tecnología.
\end{enumerate}

\section*{Referencias}

\begin{spacing}{1.0}
\small
\setlength{\parindent}{-1.27cm}
\setlength{\leftskip}{1.27cm}

Computer Security Resource Center. (s.f.). \textit{Cybersecurity}. National Institute of Standards and Technology. \url{https://csrc.nist.gov/glossary/term/cybersecurity}

Hillstone Networks. (2023, 12 de diciembre). \textit{5 diferencias entre ciberseguridad y seguridad de la información}. \url{https://www.hillstonenet.lat/blog/productos/diferencias-entre-ciberseguridad-y-seguridad-de-la-informacion/}

LISA Institute. (s.f.). \textit{Diferencia entre ciberseguridad, seguridad informática y seguridad de la información}. \url{https://www.lisainstitute.com/blogs/blog/diferencia-ciberseguridad-seguridad-informatica-seguridad-informacion}

Piranirisk. (2024, 7 de octubre). \textit{Diferencias: Seguridad de la información y ciberseguridad}. \url{https://www.piranirisk.com/es/blog/diferencias-entre-seguridad-informacion-y-ciberseguridad}

Telefónica. (2024, 25 de abril). \textit{Diferencias entre seguridad informática, seguridad de la información y ciberseguridad}. \url{https://www.telefonica.com/es/sala-comunicacion/blog/diferencias-seguridad-informatica-seguridad-informacion-ciberseguridad/}

TI Rescue. (2024, 4 de junio). \textit{Diferencias entre seguridad informática y ciberseguridad}. \url{https://tirescue.com/diferencias-entre-seguridad-informatica-y-ciberseguridad/}

\end{spacing}

\vspace{12pt}
\noindent\small Repositorio Git: \url{https://github.com/Valki11/seminario-encuesta}

\end{document}
